\begin{center}
\textbf{Explainable and Reproducible Machine Learning for ALS Prognosis: Functional Progression Classification and Survival Prediction}
\end{center}

\section*{Abstract}
\label{sec:abstract}

Amyotrophic lateral sclerosis follows a variable course, making functional decline and survival difficult to predict. This dissertation uses the PRO-ACT database in two studies of class imbalance, validation, and threshold selection. Study~I predicted rapid functional progression at three- and six-month horizons from baseline clinical features. In the six-month test set, XGBoost achieved a PR-AUC of 0.456 and detected 72 of 84 rapid progressors, with 85.7\% recall and 35.1\% precision. SHAP highlighted baseline ALSFRS-R and forced vital capacity. Study~II examined death within 24 months of symptom onset. LightGBM with random oversampling achieved a test ROC-AUC of 0.923 and detected 30 of 35 Short Survivors. Across both studies, ranking performance did not determine behaviour at the selected threshold. Outcome definition, imbalance handling, data separation, and threshold choice all affected interpretation. External, prospective validation is required before clinical use.

\section*{Keywords}
\label{sec:keywords}

Amyotrophic lateral sclerosis, functional progression, survival prognosis, class imbalance, threshold selection.
